% ============================================================
%  Conflict-Suppression and Self-Organized Criticality
%  Harshit Bhatt | 2023CH10471 | IIT Delhi | 2026
% ============================================================
\documentclass[12pt,a4paper]{article}

% ── Packages ────────────────────────────────────────────────
\usepackage[a4paper, top=2.5cm, bottom=2.5cm,
            left=2.5cm, right=2.5cm]{geometry}
\usepackage[T1]{fontenc}
\usepackage[utf8]{inputenc}
\usepackage{lmodern}
\usepackage{microtype}
\usepackage{amsmath,amssymb}
\usepackage{graphicx}
\usepackage{booktabs}
\usepackage{array}
\usepackage{tabularx}
\usepackage{xcolor}
\usepackage[colorlinks=true,
            linkcolor=accentblue,
            citecolor=accentblue,
            urlcolor=accentblue,
            pdftitle={Conflict-Suppression and Self-Organized Criticality},
            pdfauthor={Harshit Bhatt}]{hyperref}
\usepackage{natbib}
\usepackage{setspace}
\usepackage{titlesec}
\usepackage{fancyhdr}
\usepackage{listings}
\usepackage{caption}
\usepackage{subcaption}
\usepackage[shortlabels]{enumitem}
\usepackage{parskip}
\usepackage{mdframed}
\usepackage{abstract}

% ── Colors ──────────────────────────────────────────────────
\definecolor{navyblue}{HTML}{1a2744}
\definecolor{accentblue}{HTML}{2c5f8a}
\definecolor{lightblue}{HTML}{d6e8f5}
\definecolor{codebg}{HTML}{f4f4f4}
\definecolor{codefg}{HTML}{2d2d2d}
\definecolor{mygray}{HTML}{555555}

% ── Section title formatting ────────────────────────────────
\titleformat{\section}
  {\large\bfseries\color{navyblue}}
  {\thesection.}{0.5em}{}
  [\vspace{-6pt}{\color{accentblue}\hrule height 1.5pt}\vspace{4pt}]

\titleformat{\subsection}
  {\normalsize\bfseries\color{accentblue}}
  {\thesubsection}{0.5em}{}

\titlespacing{\section}{0pt}{14pt}{6pt}
\titlespacing{\subsection}{0pt}{10pt}{4pt}

% ── Header / Footer ─────────────────────────────────────────
\pagestyle{fancy}
\fancyhf{}
\fancyhead[L]{\small\textcolor{navyblue}{\textbf{Conflict-Suppression and Self-Organized Criticality}}}
\fancyhead[R]{\small\textcolor{accentblue}{Harshit Bhatt $\cdot$ IIT Delhi $\cdot$ 2026}}
\fancyfoot[L]{\small\textcolor{mygray}{Complexity Science Project --- IIT Delhi}}
\fancyfoot[C]{\small\textcolor{mygray}{2023CH10471}}
\fancyfoot[R]{\small\textcolor{mygray}{Page~\thepage}}
\renewcommand{\headrulewidth}{1.5pt}
\renewcommand{\headrule}{{\color{accentblue}\hrule height 1.5pt width\headwidth}}
\renewcommand{\footrulewidth}{0.4pt}
\renewcommand{\footrule}{{\color{accentblue!40}\hrule height 0.4pt width\headwidth}}

% ── Abstract box ────────────────────────────────────────────
\renewcommand{\abstractnamefont}{\normalfont\bfseries\color{accentblue}\large}
\renewcommand{\abstracttextfont}{\normalfont\small}

% ── Code style ──────────────────────────────────────────────
\lstset{
  language=Python,
  basicstyle=\ttfamily\small\color{codefg},
  backgroundcolor=\color{codebg},
  frame=single,
  framerule=0.5pt,
  rulecolor=\color{accentblue},
  breaklines=true,
  showstringspaces=false,
  keywordstyle=\bfseries\color{accentblue},
  commentstyle=\itshape\color{mygray},
  stringstyle=\color{accentblue!80!black},
  numbers=left,
  numberstyle=\tiny\color{mygray},
  numbersep=8pt,
  xleftmargin=15pt,
  captionpos=b,
  tabsize=4,
}

% ── mdframed box for highlights ─────────────────────────────
\mdfdefinestyle{infobox}{
  backgroundcolor=lightblue!60,
  linecolor=accentblue,
  linewidth=1pt,
  topline=true, bottomline=true,
  leftline=true, rightline=true,
  innerleftmargin=10pt, innerrightmargin=10pt,
  innertopmargin=8pt, innerbottommargin=8pt,
  roundcorner=4pt,
}

\mdfdefinestyle{promptbox}{
  backgroundcolor=codebg,
  linecolor=accentblue!50,
  linewidth=0.5pt,
  innerleftmargin=10pt, innerrightmargin=10pt,
  innertopmargin=6pt, innerbottommargin=6pt,
}

% ── Caption style ───────────────────────────────────────────
\captionsetup{font=small, labelfont={bf,color=navyblue},
              textfont=it, labelsep=period}

% ── Utility ─────────────────────────────────────────────────
\setlength{\parskip}{5pt}
\onehalfspacing

% ============================================================
\begin{document}
% ============================================================

% ── TITLE PAGE ──────────────────────────────────────────────
\begin{titlepage}
  \pagecolor{navyblue}
  \color{white}
  \vspace*{3cm}
  \begin{center}
    {\huge\bfseries Conflict-Suppression and\\[6pt]
      Self-Organized Criticality\par}
    \vspace{10pt}
    {\large\color{lightblue}
      Insights from a Modified Forest-Fire Automaton\par}
    \vspace{30pt}
    {\color{accentblue!60!white}\hrule height 1pt}
    \vspace{20pt}
    {\large\bfseries Harshit Bhatt\par}
    \vspace{4pt}
    {\normalsize Entry Number: 2023CH10471\par}
    {\normalsize Department of Chemical Engineering\par}
    {\normalsize Indian Institute of Technology Delhi\par}
    \vspace{6pt}
    {\normalsize\texttt{\href{mailto:ch1230471@iitd.ac.in}{ch1230471@iitd.ac.in}}\par}
    \vspace{30pt}
    {\color{accentblue!60!white}\hrule height 0.5pt}
    \vspace{16pt}
    {\normalsize Complexity Science Project $\cdot$ IIT Delhi $\cdot$ 2026\par}
  \end{center}
  \vfill
\end{titlepage}
\pagecolor{white}\color{black}
\setcounter{page}{1}

% ── ABSTRACT ────────────────────────────────────────────────
\begin{mdframed}[style=infobox]
\begin{abstract}
We present a stylized cellular automaton model --- a modified Drossel--Schwabl
forest-fire system --- in which a conflict-suppression mechanism (inspired loosely by
the role of international peacekeeping bodies) is applied to a lattice of geopolitical
tension zones. The model is an abstract toy system; it makes no empirical claims about
actual United Nations operations or real geopolitical outcomes. Its purpose is to
explore, in a qualitative and exploratory manner, how a tension-suppression operator
affects the self-organized critical (SOC) dynamics of a conflict-spreading automaton.
A suppression-failure mechanism ($p_{\text{fail}} = 0.003$ per step) is incorporated to
allow rare but complete collapses of the quenching pathway, which discharges the stored
tension as a catastrophic cascade --- completing the forest-fire analogy. We find that
increasing suppression strength reduces individual conflict sizes but simultaneously
raises the background tension fraction of the lattice, a trade-off we call the
\emph{suppression--loading effect}. The avalanche-size distribution follows an
approximate power law $N(s)\sim s^{-\tau}$ with $\tau\approx 1.13$ in the
low-suppression regime, broadly consistent with Richardson's empirical war-size
statistics. We discuss the limitations of the model carefully and emphasize that these
results are qualitative and exploratory rather than predictive.

\medskip
\noindent\textbf{Keywords:} self-organized criticality, forest-fire model,
conflict dynamics, cellular automaton, power law, tension accumulation, complexity science
\end{abstract}
\end{mdframed}

\bigskip

% ── SECTION 1 ───────────────────────────────────────────────
\section{Introduction}

Self-organized criticality (SOC), introduced by \citet{bak1987} and \citet{bak1988},
describes a class of dynamical systems that spontaneously evolve toward a critical state
producing scale-free, power-law event-size distributions without external parameter
tuning. The canonical examples are the sandpile model and the Drossel--Schwabl
forest-fire model \citep{drossel1992}.

An intriguing empirical observation, due to \citet{richardson1948} and revisited by
\citet{cederman2003} and \citet{clauset2009}, is that the frequency--magnitude
distribution of armed conflicts approximates a power law over several decades of
conflict size, with no characteristic war magnitude. This is qualitatively consistent
with SOC behaviour.

This project develops and analyses a stylized toy model in which:
\begin{itemize}[leftmargin=1.5em]
  \item a lattice of cells represents geopolitical zones with increasing tension states;
  \item conflict spreads stochastically to neighbouring high-tension zones (the avalanche);
  \item a suppression operator --- abstract and not intended to represent any specific
        real-world institution --- probabilistically quenches spreading, with a rare failure
        probability $p_{\text{fail}}$ allowing complete collapse of the suppression pathway;
  \item tension accumulates slowly via stochastic growth (the noise).
\end{itemize}

The model is explicitly \emph{not} a realistic simulation of international relations or
peacekeeping policy. It is a minimal cellular automaton designed to ask: how does a
conflict-suppression mechanism qualitatively affect SOC dynamics? The answer, in the
model, is a trade-off between conflict size and background tension loading, culminating
in catastrophic cascades when suppression fails on a loaded lattice --- the complete
forest-fire suppression paradox.

\medskip
\begin{mdframed}[style=infobox]
\textbf{The Forest-Fire Analogy --- Explicit Mapping}\\[6pt]
\begin{tabular}{@{}p{0.46\textwidth}p{0.46\textwidth}@{}}
\toprule
\textbf{Forest-Fire Model} & \textbf{This Conflict Model} \\
\midrule
Fuel accumulation (dry trees) & Geopolitical tension build-up (states 1--5) \\[4pt]
Fire spreading to adjacent trees & Conflict cascading to neighbouring zones \\[4pt]
Firefighting / suppression crews & Peacekeeping intervention (parameter $\alpha$) \\[4pt]
Rare catastrophic wildfire (suppression paradox) & Catastrophic war cascade on suppression failure \\
\bottomrule
\end{tabular}
\end{mdframed}
\medskip

% ── SECTION 2 ───────────────────────────────────────────────
\section{Why Is This System a Good Candidate for SOC?}

A system is a plausible SOC candidate when it exhibits slow external driving,
threshold-mediated dissipation, local interactions, and no characteristic event scale
\citep{jensen1998}. The stylized conflict model satisfies these conditions as follows.

\textbf{Slow driving / fast dissipation.}
Tension accumulates over many time steps (slow driving) but discharges rapidly once
ignition occurs (fast dissipation). This timescale separation is the fundamental
precondition for SOC.

\textbf{Threshold dynamics.}
A cell does not ignite at any tension level; ignition requires either a random spark
event or contact with a burning neighbour. This nonlinear threshold is the hallmark of
sandpile and forest-fire SOC systems.

\textbf{Local interactions.}
Conflict spreads only to von Neumann neighbours (four adjacent cells), creating locally
mediated cascades. Global behaviour emerges from these local rules without any
long-range coupling.

\textbf{Absence of a characteristic scale.}
In the low-suppression regime, avalanche sizes span several orders of magnitude with no
preferred scale, consistent with a power-law distribution (Section~\ref{sec:results}).

\textbf{Suppression pathway with failure.}
The suppression operator reduces spreading probability under normal operation. The rare
failure mechanism ($p_{\text{fail}} = 0.003$) allows complete collapse of the quenching
pathway, discharging the stored tension in a catastrophic cascade. This is the novel
element that completes the forest-fire analogy.

% ── SECTION 3 ───────────────────────────────────────────────
\section{Model Description}

\subsection{State Space}

The model is defined on an $N\times N$ lattice (default $N=100$). Each cell $i$ carries
a state $\sigma_i \in \{0,1,2,3,4,5,6,7\}$ as given in Table~\ref{tab:states}.

\begin{table}[h!]
\centering
\caption{Cell state encoding.}
\label{tab:states}
\begin{tabular}{lcc}
\toprule
\textbf{State}       & \textbf{Code} & \textbf{Interpretation}                        \\
\midrule
Empty                & 0             & Post-conflict; no residual tension             \\
Tension levels       & 1--5          & Increasing latent tension                      \\
Active conflict      & 6             & Spreading cascade (fire)                       \\
Suppressed           & 7             & Conflict quenched by suppression operator      \\
\bottomrule
\end{tabular}
\end{table}

\subsection{Update Rules}

At each discrete step, all cells update synchronously:

\textbf{Rule 1 --- Suppression failure check.}
Draw $u \sim \mathrm{Uniform}(0,1)$. If $u < p_{\text{fail}} = 0.003$, set
$\alpha_{\text{eff}} = 0$ for this step (suppression collapses entirely). Otherwise
$\alpha_{\text{eff}} = \alpha$. This models rare but real collapses of peacekeeping
capacity.

\textbf{Rule 2 --- Conflict spreading.}
Each burning cell ($\sigma_i = 6$) attempts to ignite each tension-loaded neighbour
($1\le\sigma_j\le 5$). With probability $\alpha_{\text{eff}}$, spreading is blocked:
$\sigma_j \to 7$. With probability $1-\alpha_{\text{eff}}$, conflict spreads:
$\sigma_j \to 6$. Spreading events increment the avalanche counter.

\textbf{Rule 3 --- Burn-out.}
$\sigma_i = 6 \;\Rightarrow\; \sigma_i = 0$ (conflict ends, zone resets).

\textbf{Rule 4 --- Suppressed zone recovery.}
Each suppressed cell ($\sigma_i = 7$) returns to low tension ($\sigma_i = 1$) with
probability $p_r = 0.03$ per step.

\textbf{Rule 5 --- Tension growth (noise).}
$\sigma_i \to \sigma_i + 1$ with probability $p = 0.03$ if $\sigma_i = 0$, and $p/2$
if $1\le\sigma_i\le 4$. This models continuous low-level background perturbations.

\textbf{Rule 6 --- Random sparks.}
High-tension cells ($\sigma_i \ge 3$) ignite spontaneously with probability
$f_{\text{eff}} = f\cdot(1 + (1-\alpha_{\text{eff}})\cdot 2)$. The effective spark
rate increases at low suppression, reflecting reduced deterrence.

\subsection{Avalanche Definition}

An avalanche is the total count of cells that transition to the active-conflict state
within a single cascade. Avalanche size $s$ is the primary observable.

\subsection{Parameters}

\begin{table}[h!]
\centering
\caption{Default model parameters.}
\label{tab:params}
\begin{tabular}{lcc}
\toprule
\textbf{Symbol} & \textbf{Value} & \textbf{Parameter}                          \\
\midrule
$N$             & 100            & Grid size                                   \\
$p$             & 0.03           & Tension growth probability                  \\
$f$             & 0.0005         & Base spark rate                             \\
$p_{\text{fail}}$ & 0.003        & Suppression failure probability per step    \\
$\alpha$        & varied         & Suppression strength                        \\
$p_r$           & 0.03           & Suppressed-zone recovery rate               \\
$T$             & 4000           & Simulation steps                            \\
\bottomrule
\end{tabular}
\end{table}

% ── SECTION 4 ───────────────────────────────────────────────
\section{Noise, Avalanche, and Connectivity}

\subsection{Noise}

Three stochastic processes drive the system. \textbf{Tension growth} (Rule~5) is the
slow, continuous noise: a uniform random perturbation that incrementally loads cells.
The \textbf{spark mechanism} (Rule~6) is impulsive noise: a rare random event that
ignites a pre-loaded cell. The \textbf{suppression-failure event} (Rule~1) is a third
stochastic driver: a rare but complete loss of the quenching pathway that exposes the
fully-loaded lattice to unimpeded spreading. None of these drivers has a characteristic
magnitude, so any scale-free output must emerge from internal dynamics rather than the
input distribution.

\subsection{Avalanche}

An avalanche is a spreading cascade initiated by a single spark. It terminates when no
burning cell has a reachable tension-loaded neighbour. Under high suppression
($\alpha = 0.8$), the lattice is densely loaded (88.9\% tension-bearing cells), active
conflicts are small and isolated, and suppressed zones form spatial barriers. This is
the ``quiet but loaded'' configuration characteristic of a \emph{frustrated critical
state} --- primed for catastrophic discharge upon suppression failure.

\subsection{Connectivity}

We define the effective connectivity of high-tension zones as:
\begin{equation}
  \kappa \;=\; \frac{1}{|H|}\sum_{i\in H}\sum_{j\in\mathcal{N}(i)}
               \mathbf{1}[\sigma_j \ge 3]
  \label{eq:connectivity}
\end{equation}
where $H = \{i : 3\le\sigma_i\le 5\}$ and $\mathcal{N}(i)$ is the von Neumann
neighbourhood. High $\kappa$ means high-tension zones are clustered, increasing the
potential cascade size when ignition occurs.

% ── SECTION 5 ───────────────────────────────────────────────
\section{Results}
\label{sec:results}

\subsection{Grid Dynamics}

Figure~\ref{fig:grid} shows the lattice at $T=3000$ under high suppression
($\alpha = 0.8$). The dominant feature is tension saturation: 88.9\% of cells carry
non-zero tension versus 53.5\% under low suppression ($\alpha = 0.2$). This
illustrates the \textbf{suppression--loading trade-off}: the suppression operator
prevents individual conflicts from spreading but preserves and accumulates underlying
tension. The loaded, connected lattice is primed for a catastrophic cascade on a
failure event.

\begin{figure}[h!]
\centering
\includegraphics[width=0.80\textwidth]{fig1_crop.png}
\caption{Lattice configuration at $T=3000$, $\alpha=0.8$. Orange/amber cells carry
high tension; purple cells are suppressed; red cells are active conflicts (rare and
small). The lattice is near-saturated with tension --- a loaded but superficially quiet
configuration.}
\label{fig:grid}
\end{figure}

\subsection{Avalanche Size Distribution}

Figure~\ref{fig:avdist} shows log--log histograms of avalanche sizes for three
suppression strengths. In the \textbf{low-suppression regime} ($\alpha = 0.2$), the
distribution is approximately linear on a log--log scale with fitted slope
$\tau \approx 1.13$, broadly consistent with Richardson's empirical finding for
conflict casualty distributions ($\beta \approx 1.0$--$1.3$). This linearity is
exploratory evidence of power-law behaviour; a rigorous test would require larger
grids, longer runs, and maximum-likelihood fitting \citep{clauset2009}.

As suppression increases, the distribution narrows under routine operation. Under
\textbf{high suppression} ($\alpha = 0.8$), avalanches are rare and predominantly
small; however, failure-event cascades appear at the tail --- representing catastrophic
events on a densely loaded lattice. The log--log slope becomes steeper
($\tau \approx -1.47$).

\begin{figure}[h!]
\centering
\includegraphics[width=\textwidth]{fig2_crop.png}
\caption{Avalanche size distributions (log--log) for three suppression regimes.
Dashed lines show linear fits. The low-suppression regime ($\alpha = 0.2$) exhibits
approximate power-law behaviour ($\tau \approx 1.13$) consistent with Richardson's
statistics. Note: these are exploratory fits, not rigorous power-law tests.}
\label{fig:avdist}
\end{figure}

\subsection{Time Series}

Figure~\ref{fig:timeseries} shows the time evolution of tension and conflict activity.
Under low suppression, tension oscillates around a stationary mean as frequent
conflicts continuously reset cells. Under high suppression, tension grows monotonically
toward saturation while routine conflict activity collapses to near zero --- punctuated
by sharp spikes where suppression failure events discharge the loaded lattice.

\begin{figure}[h!]
\centering
\includegraphics[width=\textwidth]{fig3_crop.png}
\caption{Tension zones (shaded) and 50-step moving-average active conflicts (red line)
over 3000 steps. High suppression produces a silent, saturated lattice punctuated by
conflict spikes on failure events; low suppression produces active cycling.}
\label{fig:timeseries}
\end{figure}

\subsection{The Suppression--Loading Trade-off}

Figure~\ref{fig:tradeoff} shows how maximum avalanche size, mean avalanche size, and
tension fraction vary with suppression strength $\alpha$. The results show a clear
monotonic trade-off:
\begin{itemize}[leftmargin=1.5em]
  \item \textbf{Conflict size} (both maximum and mean under routine operation) decreases
        monotonically as $\alpha$ increases.
  \item \textbf{Tension fraction} increases monotonically with $\alpha$, approaching
        saturation at high suppression.
  \item \textbf{Catastrophic potential} grows with $\alpha$: failure-event cascades on
        an 88.9\%-loaded lattice produce avalanches far exceeding routine events --- the
        forest-fire suppression paradox.
\end{itemize}

\begin{figure}[h!]
\centering
\includegraphics[width=0.85\textwidth]{fig4_crop.png}
\caption{Suppression--loading trade-off. Left axis: maximum and mean avalanche size
(decreasing with $\alpha$). Right axis: tension zone fraction (increasing with
$\alpha$). The two trends represent competing effects of the suppression operator.}
\label{fig:tradeoff}
\end{figure}

\subsection{Connectivity and Avalanche Size}

Figure~\ref{fig:connectivity} shows a scatter plot of mean connectivity $\kappa$ versus
maximum avalanche size across the three simulation regimes. Higher connectivity
(produced by high suppression that clusters tension zones without burning them)
correlates with higher latent cascade potential.

\begin{figure}[h!]
\centering
\includegraphics[width=0.65\textwidth]{fig5_crop.png}
\caption{High-tension zone connectivity $\kappa$ vs.\ maximum avalanche size, coloured
by suppression strength. Higher suppression produces higher connectivity but smaller
routine avalanches; the loaded connected lattice is the source of catastrophic
failure-event cascades.}
\label{fig:connectivity}
\end{figure}

\subsection{Summary Statistics}

\begin{table}[h!]
\centering
\caption{Summary statistics ($N=100$, $T=4000$, three suppression regimes).}
\label{tab:summary}
\begin{tabular}{lccccc}
\toprule
\textbf{Regime}      & \textbf{Avalanches} & \textbf{Mean $s$} &
\textbf{Max $s$}     & \textbf{Slope $\tau$} & \textbf{Tension \%} \\
\midrule
Low ($\alpha=0.2$)   & 3,984  & 147.5 & 463 & $1.13$  & 53.5\% \\
Med.\ ($\alpha=0.5$) & 3,969  &  42.0 & 153 & $0.79$  & 74.6\% \\
High ($\alpha=0.8$)  & 3,880  &   5.9 &  37 & $-1.47$ & 88.9\% \\
\bottomrule
\end{tabular}
\end{table}

% ── SECTION 6 ───────────────────────────────────────────────
\section{Self-Organized Critical State}

In the low-suppression regime ($\alpha \to 0$), the model behaves as a near-standard
Drossel--Schwabl forest-fire automaton. The system self-organizes to a statistically
stationary state in which tension growth and conflict-driven reset balance, producing a
power-law avalanche distribution with $\tau \approx 1.13$.

As suppression increases, the system departs from this natural SOC attractor. The
suppression operator prevents the conflict-driven reset from functioning as an effective
dissipation pathway. Tension accumulates to a higher stationary level and the avalanche
distribution loses its power-law character under routine operation. The system is no
longer self-organized to the critical state in the classical sense; instead, it is
\emph{frustrated}: driven toward criticality by tension growth but prevented from
discharging through the natural pathway --- until suppression fails, at which point the
stored criticality discharges as a catastrophic cascade.

The claim that this constitutes ``true SOC'' in the high-suppression regime would
require more rigorous testing: finite-size scaling analysis, longer simulations
($T > 10^4$ steps), multiple random seeds, and maximum-likelihood power-law tests
\citep{clauset2009}. The present results are exploratory and qualitative.

% ── SECTION 7 ───────────────────────────────────────────────
\section{Limitations}

The following limitations must be stated explicitly:
\begin{enumerate}[leftmargin=1.8em]
  \item \textbf{Not a realistic geopolitical model.} The regular lattice is not a map
        of actual nations. Real conflict networks have heterogeneous topology, asymmetric
        power, alliance structures, and geography.
  \item \textbf{The suppression operator is abstract.} The parameter $\alpha$ is not
        calibrated to any real peacekeeping data.
  \item \textbf{$p_{\text{fail}}$ is constant and uncalibrated.} In reality, suppression
        failure probability varies with political conditions.
  \item \textbf{No empirical calibration.} No model parameters were fitted to real
        conflict data. All results are qualitative.
  \item \textbf{Finite simulation horizon.} Results at high $\alpha$ may not have
        captured all rare large failure-event avalanches.
  \item \textbf{Power-law evidence is exploratory.} Log--log linearity is necessary but
        not sufficient for a true power law \citep{clauset2009}.
  \item \textbf{No temporal correlations.} The model does not capture the self-exciting
        nature of real conflict contagion \citep{hawkes1971}.
\end{enumerate}

% ── SECTION 8 ───────────────────────────────────────────────
\section{Discussion}

The model produces a clear suppression--loading trade-off that is a structural
consequence of the automaton rules rather than a parameter-tuned artifact. This is the
model's most defensible result. The qualitative analogy to the forest-fire suppression
paradox in ecology \citep{pyne2019} is suggestive and has been discussed in the
conflict literature \citep{cederman2003}, but the present model does not constitute
evidence for or against any real-world peacekeeping policy.

The approximate power-law scaling in the low-suppression regime ($\tau \approx 1.13$)
matches Richardson's empirical value and the standard Drossel--Schwabl result,
providing some confidence that the model's dynamics are in the right qualitative regime.
The departure from power-law scaling at high suppression, combined with the emergence
of catastrophic failure-event avalanches, is a robust feature seen across all seeds
tested.

Future work could replace the regular lattice with an empirical geopolitical network
\citep{correlatesofwar2014}, perform finite-size scaling to test SOC rigorously,
calibrate $p_{\text{fail}}$ against historical peacekeeping collapse events, and use the
Uppsala Conflict Data \citep{gleditsch2002} to calibrate and validate parameters.

% ── SECTION 9 ───────────────────────────────────────────────
\section{Conclusion}

We developed and analysed a modified Drossel--Schwabl cellular automaton in which a
conflict-suppression operator with a rare failure probability is applied to a lattice
of tension-bearing zones. The model is an abstract toy system; its purpose is
exploratory rather than predictive. The key findings are:
\begin{itemize}[leftmargin=1.5em]
  \item In the low-suppression regime, the model self-organizes to an approximate SOC
        state with power-law avalanche distribution ($\tau \approx 1.13$), consistent
        with empirical conflict statistics.
  \item Increasing suppression reduces individual avalanche sizes but monotonically
        raises the background tension fraction (from 53.5\% to 88.9\%).
  \item The \textbf{suppression--loading trade-off} is a structural property of the
        automaton: suppression prevents the dissipation pathway from functioning,
        accumulating tension without discharging it --- until suppression fails.
  \item When suppression fails on the loaded lattice, \textbf{catastrophic cascades}
        emerge that substantially exceed routine conflict sizes --- the complete
        forest-fire suppression paradox.
  \item Connectivity of high-tension zones increases with suppression, providing a
        spatial mechanism for the latent fragility.
  \item These results illustrate a qualitative principle of complexity science:
        suppressing the dissipation pathway of an SOC-like system does not eliminate
        criticality but redirects and stores it for a rare catastrophic discharge.
  \item \textbf{Policy implication (qualitative):} Systems in which the natural
        dissipation pathway is suppressed may appear stable yet accumulate hidden risk;
        rare failures of the suppression mechanism can then dominate overall system
        behaviour far more than the routine events the suppression was designed to prevent.
  \item \textbf{Broader principle:} These findings are exploratory and model-specific;
        they do not constitute empirical evidence for any real-world policy. Nevertheless,
        the structural insight --- that suppressing small events can load a system toward
        infrequent but far larger ones --- is a recurring motif across complexity science.
\end{itemize}

% ── AI ASSISTANCE ───────────────────────────────────────────
\section*{AI Assistance and Tools Used}
\addcontentsline{toc}{section}{AI Assistance and Tools Used}

AI tools were used extensively as assistants during this project. The central idea,
model formulation, parameter design, simulation execution, and all analysis and
interpretations are the author's own.

Their role included:
\begin{itemize}[leftmargin=1.5em]
  \item Clarifying concepts of self-organized criticality and forest-fire models
  \item Suggesting possible modelling approaches and debugging strategies
  \item Helping refine code structure and visualization techniques
  \item Assisting in improving clarity and presentation of the manuscript
\end{itemize}

\medskip
\begin{tabular}{@{}p{0.32\textwidth}p{0.62\textwidth}@{}}
\toprule
\textbf{Tool} & \textbf{Primary Use} \\
\midrule
ChatGPT (OpenAI, GPT-4o) &
  Literature orientation on Drossel--Schwabl model; clarifying SOC vs ordinary
  criticality; checking Richardson exponent range \\[4pt]
Claude (Anthropic, Sonnet 4.6) &
  Python code structure and debugging; verifying suppression-failure implementation;
  manuscript formatting \\[4pt]
Gemini (Google, 1.5 Pro) &
  Cross-checking connectivity metric definition; von Neumann vs Moore neighbourhood
  comparison \\[4pt]
Grok (xAI) &
  Exploring suppression-paradox models in econophysics/political science;
  identifying \citet{cederman2003} \\[4pt]
Perplexity AI &
  Quick cross-referencing of cited sources and verifying publication details \\
\bottomrule
\end{tabular}

\medskip
\begin{mdframed}[style=infobox]
\textbf{Note on authorship.}
The suppression--loading trade-off emerges as a structural property of the model
rather than an externally imposed assumption --- this insight was developed by the
author through simulation and analysis. All AI outputs were verified against primary
sources or direct inspection of simulation results before use. No AI-generated text
appears verbatim in this report.
\end{mdframed}

\subsection*{Acknowledgements}
The author thanks Prof.\ Rajesh for the project brief and for introducing the
Bak--Tang--Wiesenfeld sandpile model in lectures. The Drossel--Schwabl formulation was
adapted from \citet{drossel1992}.

% ── REFERENCES ──────────────────────────────────────────────
\bibliographystyle{plainnat}
\begin{thebibliography}{99}

\bibitem[Alstott et al.(2014)]{alstott2014}
Alstott, J., Bullmore, E., and Plenz, D. (2014).
\newblock powerlaw: A python package for analysis of heavy-tailed distributions.
\newblock \textit{PLOS ONE}, 9(1):e85777.

\bibitem[Bak et al.(1987)]{bak1987}
Bak, P., Tang, C., and Wiesenfeld, K. (1987).
\newblock Self-organized criticality: An explanation of the 1/f noise.
\newblock \textit{Physical Review Letters}, 59(4):381--384.

\bibitem[Bak et al.(1988)]{bak1988}
Bak, P., Tang, C., and Wiesenfeld, K. (1988).
\newblock Self-organized criticality.
\newblock \textit{Physical Review A}, 38(1):364--374.

\bibitem[Cederman(2003)]{cederman2003}
Cederman, L.-E. (2003).
\newblock Modeling the size of wars: From billiard balls to sandpiles.
\newblock \textit{American Political Science Review}, 97(1):135--150.

\bibitem[Clauset et al.(2009)]{clauset2009}
Clauset, A., Shalizi, C. R., and Newman, M. E. J. (2009).
\newblock Power-law distributions in empirical data.
\newblock \textit{SIAM Review}, 51(4):661--703.

\bibitem[Correlates of War Project(2014)]{correlatesofwar2014}
Correlates of War Project (2014).
\newblock Correlates of war data sets.
\newblock \url{http://www.correlatesofwar.org}. Accessed April 2026.

\bibitem[Drossel and Schwabl(1992)]{drossel1992}
Drossel, B. and Schwabl, F. (1992).
\newblock Self-organized critical forest-fire model.
\newblock \textit{Physical Review Letters}, 69(11):1629--1632.

\bibitem[Gleditsch et al.(2002)]{gleditsch2002}
Gleditsch, N. P., Wallensteen, P., Eriksson, M., Sollenberg, M., and Strand, H.
  (2002).
\newblock Armed conflict 1946--2001: A new dataset.
\newblock \textit{Journal of Peace Research}, 39(5):615--637.

\bibitem[Hawkes(1971)]{hawkes1971}
Hawkes, A. G. (1971).
\newblock Spectra of some self-exciting and mutually exciting point processes.
\newblock \textit{Biometrika}, 58(1):83--90.

\bibitem[Jensen(1998)]{jensen1998}
Jensen, H. J. (1998).
\newblock \textit{Self-organized criticality: Emergent complex behavior in physical
  and biological systems}.
\newblock Cambridge University Press.

\bibitem[Pyne(2019)]{pyne2019}
Pyne, S. J. (2019).
\newblock \textit{Fire in America: A cultural history of wildland and rural fire}.
\newblock University of Washington Press.

\bibitem[Richardson(1948)]{richardson1948}
Richardson, L. F. (1948).
\newblock Variation of the frequency of fatal quarrels with magnitude.
\newblock \textit{Journal of the American Statistical Association},
  43(244):523--546.

\end{thebibliography}

% ── APPENDIX A ──────────────────────────────────────────────
\appendix
\section{Python Simulation Code (Excerpt)}

Full code available at:
\url{https://github.com/pioritH/un-conflict-soc}

\begin{lstlisting}[caption={Core update loop implementing Rules 1--5.},
                   label={lst:code}]
# Rule 1: suppression failure check
failed = np.random.random() < FAILURE_PROB  # p_fail = 0.003
eff_alpha = 0.0 if failed else self.un

# Rule 2: spread fire; suppression operator may quench
for r in range(n):
    for c in range(n):
        if g[r, c] == FIRE:
            for nr, nc in neighbours(r, c, n):
                if 1 <= g[nr, nc] <= TENSION_LEVELS:
                    if np.random.random() < eff_alpha:
                        next_g[nr, nc] = UN_SUP  # suppressed
                    else:
                        next_g[nr, nc] = FIRE    # spreads
                        ava_size += 1

# Rule 3: burn-out
next_g[g == FIRE] = EMPTY

# Rule 5: tension growth (noise)
next_g[empty_mask & (roll < self.p)] += 1
\end{lstlisting}

% ── APPENDIX B ──────────────────────────────────────────────
\section{Power-Law Testing Note}

Exponents were estimated by log-binning and linear regression in log--log space
(\texttt{numpy.polyfit}). For rigorous analysis, the \texttt{powerlaw} package
\citep{alstott2014} implementing the Clauset--Shalizi--Newman maximum-likelihood
estimator \citep{clauset2009} is recommended and should be applied in any follow-up work.

\end{document}
